\chapter{hipótese de pesquisa}
Com base nos estudos e análises realizados até o momento, este trabalho parte 
da hipótese de que a integração do estimador de estado MUSE com o middleware 
DLS2 permitirá ao robô Go2 alcançar uma posição estimada mais precisa em terrenos 
não regulares, quando comparado ao uso de estimadores baseados exclusivamente em 
sensores proprioceptivos. Especificamente, a hipótese é que a fusão de dados 
multissensoriais, aliada a um módulo de detecção de deslizamento  e ao uso de 
observadores de atitude, será capaz de reduzir drasticamente o erro acumulado 
(drift) e evitar divergências de orientação em cenários de baixo atrito. 
Além disso, a pesquisa buscará validar essa robustez por meio de experimentos 
simulados no Gazebo, utilizando ruído sensorial controlado e ferramentas de 
análise como PlotJuggler e PCL, para mensurar quantitativamente a melhoria na 
precisão da trajetória estimada.
